As we saw in the group competition, Nepal got the chance to name a star in the constellation of Leo. The name given to that star is Sagarmatha. The coordinates of this star are $\alpha=11^{\mathrm h}35^{\mathrm m}52^{\mathrm s}$, $\delta=-4^\circ45'21''$. Given that the Autumnal equinox for the year 2024 happened on Sep 22 at 6:28 pm Nepal time, and neglecting the effects of atmospheric refraction,
\begin{enumerate}
\item Find the rise time of Sagarmatha as seen from Kathmandu on 6 October 2024. The rise time here refers to local time in Kathmandu.
\item We want to find Sagarmatha through a telescope. Since Sagarmatha is too faint, we first point to $\phi$ Leo ($\alpha=11^{\mathrm h}16^{\mathrm m}38^{\mathrm s}$, $\delta=-3^\circ38'58''$). Find the angular separation between $\phi$ Leo and Sagarmatha.
\end{enumerate}
\textit{Supplied examination data.} Kathmandu's latitude is $27.7172^\circ\mathrm{N}$.
